% appA_proofs.tex -- Supplementary Section S3: proofs of the mean first-passage time formulas (zero-based sites).
% Every displayed identity of this section is evaluated, in the form printed here, by
%   code/article/appA_proofs_check.py  ->  data/article/appA_proofs_check.json   (51 checks, 0 failed)
\section{Proofs of the mean first-passage time formulas}\label{appA_proofs:sec}

This section proves Theorems~\ref{s4_mfpt:thm-cosine}, \ref{s4_mfpt:thm-single} and~\ref{s4_mfpt:thm-pair},
Corollary~\ref{s4_mfpt:cor-resistance} and Proposition~\ref{s4_mfpt:prop-3d}. The inputs are
Proposition~\ref{s2_model:prop-pinv} and three elementary lemmas. The symbols $G_k$, $g_k$, $H_k$, $G_{kj}$,
$g_{kj}$, $\Gamma$, $O_N$ and $u_k$ are local to this section ($G_k$ is an inner sum, not the Green's function
$\Glap$). Throughout, $N\ge2$, $q\in(0,1]$, sites are
zero-based, and for $k=0,\dots,N-1$ and $m\in\mathbb Z$
\[
  \angk{k}=\frac{\pi k}{N},\qquad \ck{k}=\cos\frac{\pi k}{2N},\qquad \sk{k}=\sin\frac{\pi k}{2N},\qquad
  \thk{k}{m}=\frac{\pi k(2m+1)}{2N},
\]
so that $1-\cos\angk{k}=2\sk{k}^{2}$ and $1+\cos\angk{k}=2\ck{k}^{2}$. For $k\ge1$, $\sigk{k}=2-\cos\angk{k}>1$ and
$\phk{k}=\arcosh\sigk{k}>0$; since $\cosh\varphi=1+2\sinh^{2}(\varphi/2)$,
\begin{equation}\label{appA_proofs:eq-half-angle}
  \sinh\frac{\phk{k}}{2}=\sk{k},\qquad \coth\frac{\phk{k}}{2}=\frac{\sqrt{1+\sk{k}^{2}}}{\sk{k}} .
\end{equation}

All identities displayed below were evaluated numerically in the form printed here. The formulas, computed to
$40$ digits, were compared with exact rational solutions of $(I-\Qm)h=\one$ for the walk built from its rules
($\dd=1$: $N\le12$; $\dd=2$: $N\le9$, including all $1848$ ordered pairs with $N\le5$ at two values of $q$;
$\dd=3$: $N\le4$) and with double-precision sparse solves ($\dd=2$: $N=16,35,36,100,101$; $\dd=3$: $N=5,\dots,12$):
$51$ checks, none failed. The largest deviation is $8.5\times10^{-36}$ in the multiprecision comparisons and
$1.4\times10^{-12}$ against the sparse solves; the double-precision matrix identities (eigenvectors,
pseudo-inverses, resistances) hold to $3\times10^{-14}$.
% src: data/article/appA_proofs_check.json (n_checks, n_failed, summary, checks[].max_deviation) ; code/article/appA_proofs_check.py

\subsection{Spectrum of the walk and the cosine sum}

\begin{lemma}[Spectrum of the reflecting walk]\label{appA_proofs:lem-spectrum}
Let $\mathsf T$ be the $N\times N$ matrix with $\mathsf T(m,m\pm1)=\tfrac12$ for $0\le m,m\pm1\le N-1$,
$\mathsf T(0,0)=\mathsf T(N-1,N-1)=\tfrac12$ and zero otherwise.  Then
$\psi_k(m)=\sqrt{\alk{k}/N}\,\cos\thk{k}{m}$, $k=0,\dots,N-1$, is an orthonormal eigenbasis,
$\mathsf T\psi_k=\cos(\pi k/N)\,\psi_k$, and
$\Pm=(1-q)I+\frac{q}{\dd}\sum_{i=1}^{\dd}\mathsf T^{(i)}$, where $\mathsf T^{(i)}$ acts as $\mathsf T$ on the $i$-th coordinate.
Hence the products $\prod_i\psi_{k_i}(x_i)$ diagonalise $\Pm$ with
$1-\lambda_{\vk}=\frac{2q}{\dd}\sum_i\sk{k_i}^{2}$.
\end{lemma}

\begin{proof}
Put $u_k(m)=\cos\thk{k}{m}$ for $m\in\mathbb Z$. The addition formula gives
$u_k(m-1)+u_k(m+1)=2\cos(\angk{k})\,u_k(m)$ for every $m$. Because $\thk{k}{-1}=-\thk{k}{0}$ and
$\thk{k}{N}=2\pi k-\thk{k}{N-1}$,
\begin{equation}\label{appA_proofs:eq-bdry}
  u_k(-1)=u_k(0),\qquad u_k(N)=u_k(N-1).
\end{equation}
Hence $(\mathsf T u_k)(0)=\tfrac12[u_k(0)+u_k(1)]=\tfrac12[u_k(-1)+u_k(1)]=\cos(\angk{k})\,u_k(0)$, likewise at
$m=N-1$, and at the interior sites the eigenvalue equation is the addition formula itself. The $N$ eigenvalues
$\cos\angk{k}$ are distinct and $\mathsf T$ is symmetric, so the $u_k$ are mutually orthogonal. For $1\le k\le N-1$,
$\sum_{m=0}^{N-1}\cos^{2}\thk{k}{m}=\frac N2+\frac12\,\mathrm{Re}\sum_{m=0}^{N-1}\eu^{\iu\angk{k}(2m+1)}=\frac N2$,
because the geometric sum vanishes ($\eu^{2\iu N\angk{k}}=1$, $\eu^{2\iu\angk{k}}\ne1$); for $k=0$ the sum is $N$.
This gives the normalisation $\sqrt{\alk{k}/N}$. Finally, by the rule \eqref{s2_model:eq-P} an attempt along
coordinate $i$ has probability $q/\dd$; it changes $x_i$ by $\pm1$ with probability $\tfrac12$ each, unless the
walker would leave $\{0,\dots,N-1\}$, in which case nothing moves. That is $\mathsf T$ acting on the $i$-th
coordinate, so $\Pm=(1-q)I+\frac q\dd\sum_i\mathsf T^{(i)}$. The products are therefore eigenvectors with
$\lambda_{\vk}=1-q+\frac q\dd\sum_i\cos\angk{k_i}$, and $1-\cos\angk{k}=2\sk{k}^{2}$.
\end{proof}

As in \eqref{s2_model:eq-prop}, write $\phi_{\vk}(\vx)=\prod_i\psi_{k_i}(x_i)$. Since $\sum_i\sk{k_i}^{2}>0$ unless $\vk=\mathbf 0$, the eigenvalue
$0$ of $\Lm=I-\Pm$ is simple, and the pseudo-inverse inverts the other eigenvalues:
\begin{equation}\label{appA_proofs:eq-Lplus}
  \Lp_{\vx\vy}=\frac{\dd}{2q}\sum_{\vk\ne\mathbf 0}\frac{\phi_{\vk}(\vx)\,\phi_{\vk}(\vy)}{\sum_{i}\sk{k_i}^{2}},\qquad
  \phi_{\vk}(\vx)\,\phi_{\vk}(\vy)=\frac{1}{N^{\dd}}\prod_{i=1}^{\dd}\alk{k_i}\cos\thk{k_i}{x_i}\cos\thk{k_i}{y_i}.
\end{equation}

\begin{proof}[Proof of Theorem~\ref{s4_mfpt:thm-cosine}]
$\Pm$ is symmetric, stochastic and irreducible, so Proposition~\ref{s2_model:prop-pinv} applies with $\nn=N^{\dd}$.
Inserting \eqref{appA_proofs:eq-Lplus} into \eqref{s2_model:eq-pinv} gives \eqref{s4_mfpt:eq-cosine}; the factor
$1/q$ is explicit, so $q\,\MF_{\vo\to\va}$ does not depend on $q$. For opposite corners use
\begin{equation}\label{appA_proofs:eq-corner-cos}
  \cos\thk{k}{0}=\ck{k},\qquad \cos\thk{k}{N-1}=\cos\Bigl(\pi k-\frac{\pi k}{2N}\Bigr)=(-1)^{k}\ck{k}.
\end{equation}
The numerator of \eqref{s4_mfpt:eq-cosine} becomes $\prod_i\ck{k_i}^{2}\,[1-(-1)^{k_1+\dots+k_{\dd}}]$, which
vanishes when $k_1+\dots+k_{\dd}$ is even and equals $2\prod_i\ck{k_i}^{2}$ when it is odd. This is
\eqref{s4_mfpt:eq-cosine-cc}. For $\dd=1$, \eqref{s4_mfpt:eq-cosine} is the left-hand side of
\eqref{appA_proofs:eq-1d-sum} (because $\alk{k}=2$ and $2\sk{k}^{2}=1-\cos\angk{k}$), which equals $\hone(m,m')$ by
Lemma~\ref{appA_proofs:lem-1d} below; that lemma uses only Proposition~\ref{s2_model:prop-pinv} and
Lemma~\ref{appA_proofs:lem-spectrum}. In particular $q\,\TNd{1}=\hone(0,N-1)=N(N-1)$.
\end{proof}

\subsection{Two summation lemmas}

\begin{lemma}[Green's function of a ring]\label{appA_proofs:lem-ring}
Let $M\ge1$ be an integer and $\sigma=\cosh\varphi>1$ with $\varphi>0$.  For every integer $0\le m\le M$,
\begin{equation}\label{appA_proofs:eq-ring}
  \sum_{l=0}^{M-1}\frac{\cos(2\pi lm/M)}{\sigma-\cos(2\pi l/M)}=\frac{M\,\cosh\bigl((\tfrac M2-m)\varphi\bigr)}{\sinh\varphi\,\sinh(M\varphi/2)} ,
\end{equation}
and consequently, for every integer $0\le m\le 2N$,
\begin{equation}\label{appA_proofs:eq-ring-half}
  \sum_{j=0}^{N-1}\frac{\alk{j}\cos(\pi m j/N)}{\sigma-\cos(\pi j/N)}=\frac{2N\cosh\bigl((N-m)\varphi\bigr)}{\sinh\varphi\,\sinh N\varphi}-\frac{(-1)^{m}}{\sigma+1}.
\end{equation}
\end{lemma}

\begin{proof}
Let $A$ be the averaging operator on functions on the ring $\mathbb Z_M$, $(Af)(m)=\tfrac12[f(m-1)+f(m+1)]$. Its
eigenfunctions are $m\mapsto\eu^{2\pi\iu lm/M}$ with eigenvalues $\cos(2\pi l/M)\in[-1,1]$, so $\sigma-A$ is
invertible and
\[
  \bigl((\sigma-A)^{-1}\delta_0\bigr)(m)=\frac1M\sum_{l=0}^{M-1}\frac{\eu^{2\pi\iu lm/M}}{\sigma-\cos(2\pi l/M)},
\]
where $\delta_0$ is the indicator of the site $0$. Define
$\Gamma(m)=\cosh\bigl((\tfrac M2-m)\varphi\bigr)/\bigl(\sinh\varphi\,\sinh(M\varphi/2)\bigr)$ for $0\le m\le M$. Then
$\Gamma(0)=\Gamma(M)$, so $\Gamma$ is a function on $\mathbb Z_M$. For $1\le m\le M-1$,
$\cosh((\tfrac M2-m+1)\varphi)+\cosh((\tfrac M2-m-1)\varphi)=2\cosh\varphi\,\cosh((\tfrac M2-m)\varphi)$, that is,
$((\sigma-A)\Gamma)(m)=0$. At $m=0$ the two neighbours are $1$ and $M-1$, $\Gamma(M-1)=\Gamma(1)$, and
\[
  \sigma\,\Gamma(0)-\Gamma(1)=\frac{\cosh\varphi\,\cosh\frac{M\varphi}{2}-\cosh\bigl(\frac{M\varphi}{2}-\varphi\bigr)}
  {\sinh\varphi\,\sinh\frac{M\varphi}{2}}=1 .
\]
Hence $(\sigma-A)\Gamma=\delta_0$, so $\Gamma=(\sigma-A)^{-1}\delta_0$, and taking real parts gives
\eqref{appA_proofs:eq-ring}. For \eqref{appA_proofs:eq-ring-half} take $M=2N$: the summand
$F(l)=\cos(\pi ml/N)/(\sigma-\cos(\pi l/N))$ satisfies $F(2N-l)=F(l)$, so
$\sum_{l=0}^{2N-1}F(l)=F(0)+F(N)+2\sum_{l=1}^{N-1}F(l)=\sum_{j=0}^{N-1}\alk{j}F(j)+(-1)^{m}/(\sigma+1)$.
\end{proof}

\begin{remark}
Subtracting the term $j=0$ from \eqref{appA_proofs:eq-ring-half} and dividing by two gives
\[
  \sum_{j=1}^{N-1}\frac{\cos(\pi mj/N)}{\sigma-\cos(\pi j/N)}=\frac{N\cosh\bigl((N-m)\varphi\bigr)}{\sinh\varphi\,\sinh N\varphi}
  -\frac12\Bigl[\frac{1}{\sigma-1}+\frac{(-1)^{m}}{\sigma+1}\Bigr],
\]
which is Eq.~(E1) of \citet{Giuggioli2020}, written there with Chebyshev polynomials
($\cosh n\varphi$ and $\sinh n\varphi/\sinh\varphi$ are the Chebyshev polynomials of the first and second kind
in $\sigma$). \citet{EssamWu2009} use an equivalent summation lemma, proved by
contour integration. So Lemma~\ref{appA_proofs:lem-ring} is known; the short proof is included to keep the
section self-contained. The range $0\le m\le2N$ is sharp: the left-hand side of \eqref{appA_proofs:eq-ring-half}
is even and $2N$-periodic in $m$, the right-hand side is not, and at $m=-1$ or $m=2N+1$ the two sides differ by
exactly $4N$, because $\cosh(N+1)\varphi-\cosh(N-1)\varphi=2\sinh N\varphi\,\sinh\varphi$. For other $m$, first
reduce $m$ to $[0,2N]$.
% src: data/article/appA_proofs_check.json (labels appA_proofs:eq-E1, "appA_proofs:eq-ring-half (range)": violations 4N at sigma = 3)
\end{remark}

\begin{lemma}[One-dimensional hitting times]\label{appA_proofs:lem-1d}
For the chain $\mathsf T$ of Lemma~\ref{appA_proofs:lem-spectrum} the mean hitting time from $m$ to $m'$ is
$\hone(m,m')$ of \eqref{s4_mfpt:eq-1d}, and
\begin{equation}\label{appA_proofs:eq-1d-sum}
  \sum_{j=1}^{N-1}\alk{j}\,\frac{\cos^{2}\thk{j}{m'}-\cos\thk{j}{m}\cos\thk{j}{m'}}{1-\cos(\pi j/N)}=\hone(m,m').
\end{equation}
In particular (take $m=0$, $m'=N-1$)
\begin{equation}\label{appA_proofs:eq-cot}
  \sum_{\substack{1\le j\le N-1\\ j\ \mathrm{odd}}}\cot^{2}\frac{\pi j}{2N}=\frac{N(N-1)}{2}.
\end{equation}
\end{lemma}

\begin{proof}
Fix $m'$ and put $h(m)=\hone(m,m')$. For $m\le m'$ this is $h(m)=m'(m'+1)-m(m+1)$. It satisfies $h(m')=0$; for
$1\le m\le m'-1$, $\tfrac12[h(m-1)+h(m+1)]=h(m)-1$, because $(m-1)m+(m+1)(m+2)=2m(m+1)+2$; and at $m=0$ (if
$m'\ge1$), $(\mathsf Th)(0)=\tfrac12[h(0)+h(1)]=h(0)-1$, because $h(1)=h(0)-2$. These equations involve only sites
$\le m'$. For $m\ge m'$ the function is the same expression after the reflection $m\mapsto N-1-m$,
$m'\mapsto N-1-m'$, which maps the chain to itself, so the equations at the sites $m>m'$ hold as well. Thus
$h(m')=0$ and $((I-\mathsf T)h)(m)=1$ for all $m\ne m'$, and by the uniqueness statement of
Proposition~\ref{s2_model:prop-pinv} (the chain $\mathsf T$ is symmetric, stochastic and irreducible) $h$ is the
mean hitting time. Identity \eqref{appA_proofs:eq-1d-sum} is \eqref{s2_model:eq-pinv} for $\mathsf T$, with
$\nn=N$ and $(I-\mathsf T)^{+}_{mm'}=\sum_{j\ge1}\psi_j(m)\psi_j(m')/(1-\cos\angk{j})$ from
Lemma~\ref{appA_proofs:lem-spectrum}. For \eqref{appA_proofs:eq-cot} take $m=0$, $m'=N-1$ and use
\eqref{appA_proofs:eq-corner-cos}: the summand vanishes for even $j$ and equals $2\ck{j}^{2}/\sk{j}^{2}$ for odd
$j$, and $\hone(0,N-1)=N(N-1)$.
\end{proof}

\subsection{Proof of Theorem~\ref{s4_mfpt:thm-single}}

\emph{Step 1 (splitting off the row $k=0$).} For $\dd=2$, $\sk{k}^{2}+\sk{j}^{2}=\tfrac12(2-\cos\angk{k}-\cos\angk{j})$,
and \eqref{s4_mfpt:eq-cosine-cc} reads
$\tfrac12\,q\,\TN=\sum_{(k,j)\ne(0,0)}\alk{k}\alk{j}\ck{k}^{2}\ck{j}^{2}[1-(-1)^{k+j}]/(2-\cos\angk{k}-\cos\angk{j})$.
In the row $k=0$ the denominator is $1-\cos\angk{j}$ and $\alk{0}\ck{0}^{2}=1$; in a row $k\ge1$ it is
$\sigk{k}-\cos\angk{j}$, $\alk{k}=2$, and $j$ runs over $0,\dots,N-1$. Therefore
\begin{equation}\label{appA_proofs:eq-split}
  \frac{q\,\TN}{2}=\sum_{j=1}^{N-1}\frac{\alk{j}\ck{j}^{2}\,[1-(-1)^{j}]}{1-\cos\angk{j}}+\sum_{k=1}^{N-1}2\ck{k}^{2}\,G_k,
  \qquad
  G_k:=\sum_{j=0}^{N-1}\frac{\alk{j}\ck{j}^{2}\,[1-(-1)^{k+j}]}{\sigk{k}-\cos\angk{j}} .
\end{equation}
By \eqref{appA_proofs:eq-1d-sum} with $m=0$, $m'=N-1$ the first sum equals $\hone(0,N-1)=N(N-1)$.

\emph{Step 2 (the inner sum).} Write $\ck{j}^{2}=\tfrac12[(1+\sigk{k})-(\sigk{k}-\cos\angk{j})]$ and
$(-1)^{j}=\cos(N\angk{j})$. Then
\[
  G_k=\frac{1+\sigk{k}}{2}\sum_{j=0}^{N-1}\frac{\alk{j}\,[1-(-1)^{k}\cos(N\angk{j})]}{\sigk{k}-\cos\angk{j}}
  -\frac12\sum_{j=0}^{N-1}\alk{j}\,[1-(-1)^{k+j}] .
\]
By \eqref{appA_proofs:eq-ring-half} with $m=0$ and with $m=N$,
\begin{align*}
  \sum_{j=0}^{N-1}\frac{\alk{j}}{\sigk{k}-\cos\angk{j}}&=\frac{2N\cosh N\phk{k}}{\sinh\phk{k}\,\sinh N\phk{k}}-\frac{1}{\sigk{k}+1},\\
  \sum_{j=0}^{N-1}\frac{\alk{j}\cos(N\angk{j})}{\sigk{k}-\cos\angk{j}}&=\frac{2N}{\sinh\phk{k}\,\sinh N\phk{k}}-\frac{(-1)^{N}}{\sigk{k}+1},
\end{align*}
and $\sum_{j=0}^{N-1}\alk{j}=2N-1$, $\sum_{j=0}^{N-1}\alk{j}(-1)^{j}=-(-1)^{N}$ (from
$\sum_{j=0}^{2N-1}(-1)^{j}=0$ and the symmetry $j\mapsto2N-j$). The terms containing $(-1)^{k+N}$ cancel, the
constants combine to $-N$, and
\begin{equation}\label{appA_proofs:eq-Gk}
  G_k=\frac{N\,(1+\sigk{k})\,[\cosh N\phk{k}-(-1)^{k}]}{\sinh\phk{k}\,\sinh N\phk{k}}-N .
\end{equation}

\emph{Step 3 (form (a)).} Since $\cosh(N-1)\varphi=\cosh N\varphi\cosh\varphi-\sinh N\varphi\sinh\varphi$ and
$\cosh\phk{k}=\sigk{k}$,
\begin{equation}\label{appA_proofs:eq-Bk}
  \Bk{k}=(1+\sigk{k})\,[\cosh N\phk{k}-(-1)^{k}]-\sinh\phk{k}\,\sinh N\phk{k},
  \qquad G_k=\frac{N\,\Bk{k}}{\sinh\phk{k}\,\sinh N\phk{k}} ,
\end{equation}
the second relation being \eqref{appA_proofs:eq-Gk} rewritten. Inserting it and $N(N-1)$ into \eqref{appA_proofs:eq-split} gives \eqref{s4_mfpt:eq-single-a}. The rows
$k\ge1$ were summed over all $j$, including $j=0$; no cancellation between different sectors of the double sum is
needed.

\emph{Step 4 (form (b)).} Because $(1+\cosh\varphi)/\sinh\varphi=\coth(\varphi/2)$,
$(\cosh N\varphi-1)/\sinh N\varphi=\tanh(N\varphi/2)$ and $(\cosh N\varphi+1)/\sinh N\varphi=\coth(N\varphi/2)$,
\eqref{appA_proofs:eq-Gk} reads
\begin{equation}\label{appA_proofs:eq-Gk-coth}
  G_k=N\Bigl[\coth\frac{\phk{k}}{2}\,g_k-1\Bigr],\qquad
  g_k=\coth\frac{N\phk{k}}{2}\ (k\ \text{odd}),\qquad g_k=\tanh\frac{N\phk{k}}{2}\ (k\ \text{even}).
\end{equation}
Moreover
\begin{equation}\label{appA_proofs:eq-csum}
  \sum_{k=1}^{N-1}\ck{k}^{2}=\frac{N-1}{2}+\frac12\sum_{k=1}^{N-1}\cos\angk{k}=\frac{N-1}{2},
\end{equation}
since the terms $k$ and $N-k$ of the cosine sum cancel. Hence, by \eqref{appA_proofs:eq-split},
\[
  q\,\TN=2N(N-1)+4N\sum_{k=1}^{N-1}\ck{k}^{2}\coth\frac{\phk{k}}{2}\,g_k-4N\cdot\frac{N-1}{2}
        =4N\sum_{k=1}^{N-1}\ck{k}^{2}\coth\frac{\phk{k}}{2}\,g_k ,
\]
which is \eqref{s4_mfpt:eq-single-S} by \eqref{appA_proofs:eq-half-angle}; note $N\phk{k}/2=N\arsinh\sk{k}$.

\emph{Step 5 (form (c)).} The summand of \eqref{s4_mfpt:eq-cosine-cc} is symmetric in $(k,j)$, and each pair with
$k+j$ odd has exactly one odd index. With $1/(\sk{k}^{2}+\sk{j}^{2})=2/(\sigk{k}-\cos\angk{j})$,
\begin{equation}\label{appA_proofs:eq-odd-even}
  q\,\TN=16\sum_{\substack{1\le k\le N-1\\ k\ \mathrm{odd}}}\ck{k}^{2}\,H_k,\qquad
  H_k:=\sum_{\substack{0\le j\le N-1\\ j\ \mathrm{even}}}\frac{\alk{j}\ck{j}^{2}}{\sigk{k}-\cos\angk{j}} .
\end{equation}
For $j=2l$ we have $\angk{j}=2\pi l/N$ and $\ck{j}^{2}=\cos^{2}(\pi l/N)$. The function
$l\mapsto\cos^{2}(\pi l/N)/(\sigk{k}-\cos(2\pi l/N))$ on $\mathbb Z_N$ is invariant under $l\mapsto N-l$ and
vanishes at $l=N/2$ when $N$ is even; hence $H_k$ equals its sum over all $l\in\mathbb Z_N$, for both parities of
$N$. Writing $\cos^{2}(\pi l/N)=\tfrac12[(1+\sigk{k})-(\sigk{k}-\cos(2\pi l/N))]$ and using
\eqref{appA_proofs:eq-ring} with $M=N$, $m=0$,
\begin{equation}\label{appA_proofs:eq-Hk}
  H_k=\frac{1+\sigk{k}}{2}\cdot\frac{N\cosh(N\phk{k}/2)}{\sinh\phk{k}\,\sinh(N\phk{k}/2)}-\frac N2
     =\frac N2\Bigl[\coth\frac{\phk{k}}{2}\coth\frac{N\phk{k}}{2}-1\Bigr].
\end{equation}
Next,
\begin{equation}\label{appA_proofs:eq-csum-odd}
  \sum_{\substack{1\le k\le N-1\\ k\ \mathrm{odd}}}\ck{k}^{2}=\frac N4\qquad\text{for every }N\ge2 .
\end{equation}
Indeed the left-hand side is half the number of odd $k<N$ plus $\tfrac12\,O_N$, where $O_N$ denotes the sum of
$\cos\angk{k}$ over the odd $k<N$. For even $N$ the number is $N/2$ and $O_N=0$, since $k\mapsto N-k$ maps odd
indices to odd indices and changes the sign of the cosine. For odd $N$ the number is $(N-1)/2$ and
$O_N=\tfrac12$. To see this, note that for odd $N$ the alternating sum
\[
  \sum_{k=0}^{N-1}(-1)^{k}\eu^{\iu\angk{k}}=\frac{1-(-1)^{N}\eu^{\iu\pi}}{1+\eu^{\iu\pi/N}}
\]
vanishes, so the sum of $\cos\angk{k}$ over the even $k$ (which includes the term $k=0$, equal to $1$) equals
$O_N$; on the other hand $k\mapsto N-k$ maps the even indices $k\ge2$ onto the odd ones and changes the sign of
the cosine, so the same sum equals $1-O_N$. In both cases the left-hand side is $N/4$. From
\eqref{appA_proofs:eq-odd-even}--\eqref{appA_proofs:eq-csum-odd},
\[
  q\,\TN=8N\sum_{k\ \mathrm{odd}}\ck{k}^{2}\Bigl[\coth\frac{\phk{k}}{2}\coth\frac{N\phk{k}}{2}-1\Bigr]=8N\,S_{\mathrm{odd}}-2N^{2}.
\]
Form (b) is $q\,\TN=4N\,(S_{\mathrm{odd}}+S_{\mathrm{even}})$. Comparing the two gives
$S_{\mathrm{odd}}-S_{\mathrm{even}}=N/2$ and then $q\,\TN=2N^{2}+8N\,S_{\mathrm{even}}$, which completes
\eqref{s4_mfpt:eq-half}. \qed

\begin{remark}
Form (a) contains $\cosh N\phk{k}$ and $\sinh N\phk{k}$, which are of size $\eu^{N\phk{k}}/2$. Since
$\phk{N-1}\to\arcosh3=1.7627$, the exponent $N\phk{N-1}$ passes the logarithm of the largest double-precision
number, $709.78$, between $N=402$ ($708.62$) and $N=403$ ($710.38$). Accordingly a direct double-precision
evaluation of \eqref{s4_mfpt:eq-single-a} is finite for every $N\le402$ (at $N=402$ it agrees with form (b) to
$5.7\times10^{-13}$) and returns no finite value at any size tested from $N=403$ on.
% src: data/article/appA_proofs_check.json (float64_overflow)
Forms (b) and (c) contain only $\tanh$ and $\coth$ and have no such limit.
\end{remark}

\subsection{Arbitrary start and target, resistance, and three dimensions}

\begin{proof}[Proof of Theorem~\ref{s4_mfpt:thm-pair}]
For $\dd=2$, \eqref{s4_mfpt:eq-cosine} reads
$q\,\MF_{\vo\to\va}=\sum_{(k,j)\ne(0,0)}2\alk{k}\alk{j}\,[\mathcal C_k(a_1,a_1)\mathcal C_j(a_2,a_2)-\mathcal C_k(o_1,a_1)\mathcal C_j(o_2,a_2)]/(2-\cos\angk{k}-\cos\angk{j})$.
In the row $k=0$, $\alk{0}=1$, $\mathcal C_0\equiv1$ and the denominator is $1-\cos\angk{j}$, so by
\eqref{appA_proofs:eq-1d-sum} this row contributes $2\,\hone(o_2,a_2)$. In the rows $k\ge1$ the sums over $j$ have
the form $\sum_{j=0}^{N-1}\alk{j}\cos\thk{j}{m}\cos\thk{j}{m'}/(\sigk{k}-\cos\angk{j})$. Now
\begin{equation}\label{appA_proofs:eq-prod}
  \cos\thk{j}{m}\cos\thk{j}{m'}=\tfrac12\bigl[\cos\bigl(|m-m'|\,\angk{j}\bigr)+\cos\bigl((m+m'+1)\,\angk{j}\bigr)\bigr],
\end{equation}
where $m_1=|m-m'|\in[0,N-1]$ and $m_2=m+m'+1\in[1,2N-1]$ have opposite parity. Both lie in the range of
\eqref{appA_proofs:eq-ring-half}; applying it to the two terms, the contributions $-(-1)^{m_1}/(\sigk{k}+1)$ and
$-(-1)^{m_2}/(\sigk{k}+1)$ cancel and
\begin{equation}\label{appA_proofs:eq-inner-pair}
  \sum_{j=0}^{N-1}\frac{\alk{j}\cos\thk{j}{m}\cos\thk{j}{m'}}{\sigk{k}-\cos\angk{j}}=N\,\Psi_k(m,m'),
\end{equation}
with $\Psi_k$ as in \eqref{s4_mfpt:eq-Psi}. Since $\alk{k}=2$ for $k\ge1$, this gives \eqref{s4_mfpt:eq-pair}.
\end{proof}

For $\vo=(0,0)$ and $\va=(N-1,N-1)$, \eqref{appA_proofs:eq-corner-cos} gives $\mathcal C_k(a_1,a_1)=\ck{k}^{2}$ and
$\mathcal C_k(o_1,a_1)=(-1)^{k}\ck{k}^{2}$, the numerators of $\Psi_k(N-1,N-1)$ and $\Psi_k(0,N-1)$ are
$\cosh N\phk{k}+\cosh(N-1)\phk{k}$ and $\cosh\phk{k}+1$, and \eqref{s4_mfpt:eq-pair} reduces to
\eqref{s4_mfpt:eq-single-a}. For numerical work one writes
$\cosh((N-m)\varphi)/\sinh N\varphi=(\eu^{-m\varphi}+\eu^{-(2N-m)\varphi})/(1-\eu^{-2N\varphi})$, valid for
$0\le m\le2N$, which removes the overflow.

\begin{proof}[Proof of Corollary~\ref{s4_mfpt:cor-resistance}]
Let $\mathcal L=D-A$ be the Laplacian of the $N\times N$ grid graph, with $D$ its degree matrix and $A$ its
adjacency matrix. By \eqref{s2_model:eq-P}, $\Pm(\vx,\vy)=q/4$ for neighbours and
$1-\Pm(\vx,\vx)=\tfrac q4\deg(\vx)$, so
\begin{equation}\label{appA_proofs:eq-laplacian}
  I-\Pm=\frac q4\,\mathcal L,\qquad \Lp=\frac4q\,\mathcal L^{+}.
\end{equation}
Proposition~\ref{s2_model:prop-pinv}, applied to both directions, gives
$\MF_{\vo\to\va}+\MF_{\va\to\vo}=\nn\,(\Lp_{\va\va}+\Lp_{\vo\vo}-2\Lp_{\vo\va})
=\frac{4N^{2}}{q}\,(\ve_{\va}-\ve_{\vo})^{\T}\mathcal L^{+}(\ve_{\va}-\ve_{\vo})$, with $\ve_{\vx}$ the unit vector of
site $\vx$, and the last quadratic form is the effective resistance $\Reff(\vo,\va)$ \citep{KleinRandic1993}. This is
\eqref{s2_model:eq-resistance} for $\dd=2$, the commute-time identity of \citet{Chandra1989}; see also
\citet{Tetali1991}. The point reflection $\vx\mapsto(N-1-x_1,N-1-x_2)$ is an automorphism of the walk that
exchanges the two corners, so for them $\MF_{\vo\to\va}=\MF_{\va\to\vo}$ and $q\,\TN=2N^{2}\Reff_N$. The two
expressions for $\Reff_N$ follow from \eqref{s4_mfpt:eq-half}. For example $q\,T_{4}=416/7$ gives $\Reff_4=13/7$,
the value that Example~4 of \citet{Wu2004} gives for unit resistors.
% src: data/article/appA_proofs_check.json (exact_corner_values_d2; check "exact: q T_4/(2*16) = 13/7")
\end{proof}

\begin{proof}[Proof of Proposition~\ref{s4_mfpt:prop-3d}]
Put $\sigma_{kj}=3-\cos\angk{k}-\cos\angk{j}=1+2\sk{k}^{2}+2\sk{j}^{2}$, so that
$\sk{k}^{2}+\sk{j}^{2}+\sk{l}^{2}=\tfrac12(\sigma_{kj}-\cos\angk{l})$ and, for $(k,j)\ne(0,0)$,
$\sigma_{kj}=\cosh\varphi_{kj}>1$. Summing \eqref{s4_mfpt:eq-cosine-cc} for $\dd=3$ over the third index first,
\begin{equation}\label{appA_proofs:eq-Gkj}
  q\,\TNd{3}=3\sum_{k,j=0}^{N-1}\alk{k}\alk{j}\ck{k}^{2}\ck{j}^{2}\,G_{kj},\qquad
  G_{kj}:=\sum_{l}\frac{\alk{l}\ck{l}^{2}\,[1-(-1)^{k+j+l}]}{\sigma_{kj}-\cos\angk{l}},
\end{equation}
where $l$ runs over $0,\dots,N-1$, except for $(k,j)=(0,0)$, where $l\ge1$. As in Step~1, $G_{00}=N(N-1)$ by
\eqref{appA_proofs:eq-1d-sum}. For $(k,j)\ne(0,0)$, Steps~2 and~4 apply verbatim with $\sigk{k}$, $\phk{k}$ and
$(-1)^{k}$ replaced by $\sigma_{kj}$, $\varphi_{kj}$ and $(-1)^{k+j}$:
$G_{kj}=N[\coth(\varphi_{kj}/2)\,g_{kj}-1]$. Finally
$\sum_{k=0}^{N-1}\alk{k}\ck{k}^{2}=1+2\cdot\frac{N-1}{2}=N$ by \eqref{appA_proofs:eq-csum}, so
$\sum_{(k,j)\ne(0,0)}\alk{k}\alk{j}\ck{k}^{2}\ck{j}^{2}=N^{2}-1$ and
\[
  q\,\TNd{3}=3\Bigl[N(N-1)+N\!\!\sum_{(k,j)\ne(0,0)}\!\!\alk{k}\alk{j}\ck{k}^{2}\ck{j}^{2}\coth\frac{\varphi_{kj}}{2}\,g_{kj}-N(N^{2}-1)\Bigr],
\]
which is \eqref{s4_mfpt:eq-3d}, because $N(N-1)-N(N^{2}-1)=-N^{2}(N-1)$.
\end{proof}
